\documentclass[12pt]{article}
\usepackage{geometry}
\geometry{letterpaper, margin=1in}
\usepackage{hyperref}
\usepackage[usenames,dvipsnames]{xcolor}
\usepackage{natbib}
\usepackage{microtype}
\bibliographystyle{apalike}

\definecolor{accent}{RGB}{26, 60, 110}   % matches resume navy

\begin{document}

%----------HEADER----------
\begin{center}
  {\LARGE\textbf{\textcolor{accent}{SITONG LI}}} \\[4pt]
  \small Minneapolis, MN
  \;\textbullet\;
  \href{mailto:ruc.sitongli@gmail.com}{ruc.sitongli@gmail.com}
  \;\textbullet\;
  \href{https://linkedin.com/in/YOURHANDLE}{LinkedIn}   % <-- fill in
  \;\textbullet\;
  \href{https://github.com/YOURHANDLE}{GitHub}          % <-- fill in
  \;\textbullet\;
  (952) 201-4026
\end{center}

\vspace{1em}

\noindent \today

\vspace{1em}

% ============================================================
% FILL IN: e.g., \noindent Dear Professor [Name],
% ============================================================
\noindent Dear Hiring Committee,

\vspace{1em}

I am applying for the predoctoral fellowship and am
strongly committed to pursuing a Ph.D.\ in Economics.  My primary
research interests lie in political economy---in both democratic and
non-democratic settings.  In democracies, I want to understand how
political dynamics, including election cycles and populist rhetoric,
shape regulatory behavior, media incentives, and firms' decisions,
with downstream effects on fiscal policy and economic outcomes.  Prior
work, including \cite{kempf2021partisan} and \cite{fos2022political},
documents how polarization is reshaping financial markets, and
\cite{cassidy2025partisan} shows that firms adjust their public
communication in response.  Given the broader rise of populism across
democratic countries, I see this as a timely and consequential
research agenda.

I am equally interested in non-democratic settings.  In authoritarian
regimes such as contemporary China, I want to study how information
control---censorship and propaganda---distorts the information
environment for markets and intermediaries, with consequences for
firms, investors, and ordinary citizens.  \cite{lorentzen2014china}
documents the strategic logic behind censorship in China, and
\cite{chen2019impact} shows that it can suppress information-seeking.
Building on these insights, I am interested in how censorship shapes
investment behavior and capital allocation.

My interest in non-democratic political economy is partly personal.  I
grew up in China, where politics is not something you observe from a
distance---it is something you live with every day.  When I began my
career doing fixed-income research at Harvest Fund, I assumed that
rigorous quantitative analysis would be sufficient.  But while working
on municipal bonds, I found that no purely technical model can
substitute for understanding local fiscal capacity, officials' policy
priorities, or the credibility of implicit government guarantees.
Bond prices in China depend heavily on the commitments of local
officials; when an official is removed or investigated under an
anti-corruption campaign, policy can shift and the economic
environment can change overnight.  This information is often censored
or strategically withheld, yet it is decisive for pricing.  That
experience was the origin of my interest in non-democratic
institutions and the economics of information control.

After coming to the United States and receiving research training with
Professors Cassidy and Zhang, I found that the same themes arise even
in democratic settings.  From Prof.\ Cassidy's project, I learned how
congressional committee dynamics shape public discourse on fiscal
policy and how legislators bring political pressure to bear on federal
agencies.  From Prof.\ Zhang's projects, I explored how the EPA
responds to congressional information requests and how firms' climate
risk disclosures vary with regulatory and political pressure.  These
experiences deepened my conviction that understanding political
institutions is essential for empirical work in finance and
economics---which is why I am applying for this predoc.

My research training has also given me broad methodological experience.
I have worked with both classical NLP approaches---bag-of-words
models, named-entity recognition, and rule-based classifiers---and
modern text-as-data methods in political economy, including MNIR
\citep{taddy2013multinomial} as applied in
\cite{gentzkow2019measuring}, BERT-based classifiers, and LLM-based
extraction workflows.  I care as much about measurement quality as I
do about statistical power: with the rise of LLMs, extracting
constructs from text has become easier, but measurement error can
also grow if workflows are not transparent or reproducible.  I pay
careful attention to practical steps that improve robustness---
controlling temperature, using cross-validation, and checking results
against alternative NLP specifications.  I read widely in the
textual-analysis literature and have discussed methods directly with
authors to understand not only how they manage measurement error, but
also how they communicate those choices to editors and discussants.
For instance, I suggested that our team follow \cite{hampole2025artificial}
and run a local LLM as a robustness check for our GPT-4-extracted
measures.

As a predoctoral researcher, I aim to contribute in three concrete
ways.  First, I can build and maintain scalable, reproducible data
pipelines for large datasets.  Second, I can develop measurement
strategies for narratives and information environments using NLP and
LLMs, with close attention to robustness and interpretability.  Third,
I can support credible empirical designs---clean panel construction,
event-study and difference-in-differences frameworks, and
replication-ready code.

I would be grateful for the opportunity to work with your group and to
deepen my training in economics research.  I am confident that my
background---combining firsthand experience with non-democratic
institutions, substantive interest in censorship and information
economics, and strong technical skills in data and text
analysis---would allow me to contribute meaningfully to your research
agenda.

\vspace{1em}

\noindent Sincerely,\\[0.5em]
\textbf{Sitong Li}

\bibliography{refs}

\end{document}
